\scanpage{202}
\mspage{58}
est la catégorie disconnexe $\unsure{\varepsilon|2}=\{\bullet\ \unsure{\bullet|{,}}\}$.

J'avais espéré un moment que les foncteurs $W_0$-connexes $I'\xrightarrow{u}\unsure{I|I'}$ seraient ceux qui se factorisent en $I'\overset{u_0}{\mathbin{\unsure{-|\to}}}I''\overset{v}{\mathbin{-}}I$, avec \deleted{\unsure{$I_0$}} $u_0$ une équivalence, et $v$ une $W_0$-équivalence universelle. Mais l'exemple précédent le met en défaut, avec $u$ \inserted{même} \deleted{\unsure{une}} $W_\infty$-connexe. En effet, $I''$ serait \deleted{\unsure{un}} déduit d'un ens. préordonné, et $v$ se factoriserait par l'ens ordonné correspondant, savoir $I'$, donc serait de la forme

\[
 I''\longrightarrow I'\xrightarrow{u}I,
\]

où $I''\mathbin{-}I'$ est un passage au quotient, et l'image inverse de $\unsure{\mathbb{\Delta}|\Delta}^1\overset{\partial}{\mathbin{-}}I$ serait encore disconnexe (un $\pi_0$ de cardinal $2$). Donc c'est vraiment sans espoir.

Cependant

\uline{Corollaire} 4.3. Supposons que $u$ soit $W_0$-lisse, ou $W_0$-propre, ou une $W_0$-fibration faible. Alors $u$ est $W_0$-connexe sss c'est une $W_0$-équiv. universelle, i.e. ssi a) ses fibres sont $0$-connexes et b) $I'\times_I\mathbb{\Delta}^1$ $0$-connexe pour $\forall\ \mathbb{\Delta}^{\unsure{1}}\mathbin{-}I$ (correspondant à une flèche non identique).
% archive sheet number: 202
